\section{Teoretické statistické vlastnosti PFN}
\label{sec:pfn-teorie}
V této sekci uvedeme dosud známé statistické vlastnosti PFN, které publikoval Nagler v článku (Nagler, 2023).
V okamžiku, kdy již máme natrénovaný model, dává smysl prozkoumat jeho statistické vlastnosti a ptat se, proč a jak fungují jeho prediktivní a aproximační schopnosti.
Muller ve své práci ukázal empirická pozorování, která ale nepodpořil teoretickým zdůvodněním.
To se avšak podařilo Naglerovi, který dokázal vybudovat teoretický základ pro PFN a řádně odůvodnit, proč a jak architektura PFN dává dobré výsledky při Bayesovské inference.
Nagler ve své práci používá určitý formalismus pro zavedení celého problému, na kterém pak buduje samotnou teorii.
My tento formalismus v této práci zachováme.

\subsection{Statistický model}
\label{subsec:statisticky-model}
Uvažujme klasifikační úlohu, která se skládá ze tříd $Y \in \mathcal{Y}$ a příznaků $X \in \mathcal{X} \subseteq \mathbb{R}^d$.
Předpokládejme dále, že máme iid data $D_n = (Y_i, X_i)_{i=1}^n$ pocházející z nějakého teoretického, avšak neznámého rozdělení $p_0$.
Naším cílem je predikovat toto podmíněné rozdělení na základě dostupných dat: $$p_0(y \mid x) = P(Y = y \mid X = x).$$
Z hlediska Baysevského neparametrického odhadování můžeme na $p_0$ nahlížet jako na realizaci náhodného parametru $p \in \mathcal{P}$, kde $\mathcal{P}$ je prostor podmíněných pravděpodobnostních rozdělení nad $\mathcal{Y} \times \mathcal{X}$.
Tento parametr má rozdělení $\Pi$, které budeme říkat prior.
Prior bude vyjadřovat naše přesvědčení o tom, jaké modely $p$ jsou pravděpodobnější pro daný problém ještě předtím než uvidíme data.
Formálně celý proces, ze kterého pocházejí data lze zapsat takto:
\begin{enumerate}
    \item Vygeneruj $p \sim \Pi$.
    \item Vygeneruj iid vzorky $D_n = (Y_i, X_i)_{i=1}^n$ a jeden dodatečný pár $(Y, X)$ z modelu $p$.
\label{generace_dat}
\end{enumerate}

Tento mechanismus definuje sdílené rozdělení nad $(D_n \cup (Y, X), p)$ pro každé $n$. 
Skutečné rozdělení $p_0$ je ta realizace $p$, ze které skutečně pocházejí naše data, zpravidla je $p_0$ neznámá.
\par   

Pro každé $n$ poskytuje výše popsaný statistický model dobře definované sdružené rozdělení n-tice $(D_n \cup (Y, X), p)$. 
Na základě tohoto modelu lze $p_0(y \mid x)$ aproximovat pomocí posterior predictive distribution (PPD):
$$\pi(y \mid x, D_n) = P(Y = y \mid X = x, D_n),$$
kterou jsme již zmiňovali v předchozí sekci, uveďme její definici ještě jednou i zde a se zařazením do kontextu.
Tím vzniká celá rodina PPD indexovaná $n$. 
Pokud prior $\Pi$ faktorizuje nezávisle pro marginální složky $p(y \mid x)$ a $p(x)$, tj. rozdělení příznaků a rozdělení tříd jsou v prioru nezávislé, pak lze PPD dále zapsat jako:
\begin{equation}
    \pi(y \mid x, D_n) = \int p(y \mid x) \, d\Pi(p \mid D_n),
\label{def:PPD}
\end{equation}

kde $\Pi(\cdot \mid D_n)$ je tzv. \textit{posterior}, což je podmíněné rozdělení $p$ dané pozorovanými daty $D_n$, získaná Bayesovým pravidlem. 
Tento integrál říká v zásadě relativně jednoduchou věc: průměrujeme $$p(y \mid x)$$ přes všechna možná $p$ vážená jejich posteriorní pravděpodobností, pak modely $p$, které jsou konzistentní s $D_n$ dostanou větší váhu.
Podmínka faktorizace prioru je spíše technická.
Platí–li podmínka, pak když jsme viděli data $D_n$, říkají nám něco o $p(y \mid x)$, ale $x$–ová souřadnice testovacího bodu sáma o sobě nic neříká o $p(y \mid x)$.
Pokud by totiž ta podmínka neplatila, potom by i testovací souřadnice $x$ byla informativní pro $p(y \mid x)$, což pro náš případ nedává smysl.
PPD $\pi(y \mid x, D_n)$ je tedy \textit{posteriorní střední hodnota} podmíněných rozdělení $p(y \mid x)$, vážená tím, jak jsou jednotlivé modely $p$ konzistentní s $D_n$.
\par   

PFN je numerická aproximace celé rodiny PPD $\{\pi(\cdot \mid \cdot, D_n),\, n \in \mathbb{N}\}$.
Základním teoretickým poznatkem je, že PPD sama o sobě je pro každé $n$ optimálním prediktorem v následujícím smyslu. 
Definujme třídu všech podmíněných pravděpodobnostních funkcí takto:
$$\mathcal{Q} = \left\{ q : (\mathcal{Y} \times \mathcal{X})^{n+1} \to [0,1] \;\Big|\; \sum_{y \in \mathcal{Y}} q(y \mid \cdot, \cdot) = 1 \right\}.$$
Každé $q \in \mathcal{Q}$ je tedy funkce, která dostane $n$ trénovacích párů $(Y_i, X_i)$ a jeden testovací vstup $\mathbf{X}$, a vrátí distribuci přes $\mathcal{Y}$.
\begin{theorem}
    $\pi$ z rovnice \ref{def:PPD} splňuje:
    \begin{equation}
            \pi = \arg\max_{q \in \mathcal{Q}} \, \mathbb{E}_\Pi[\log q(Y \mid X, D_n)],
    \end{equation}
    kde $\mathbb{E}_\Pi$ je střední hodnota přes $(Y, X) \cup D_n$ generované dle mechanismu v sekci 2.1.
\label{theorem:prior_je_max}
\end{theorem}
Jinými slovy PPD $\pi$ maximalizuje očekávanou podmíněnou log-likelihood přes všechny možné prediktory z $\mathcal{Q}$. 
To znamená, že to je nejlepší možný prediktor za daného prioru $\Pi$.

Dále navíc maximalizaci výrazu $\mathbb{E}_\Pi[\log q(Y \mid X, D_n)]$ lze ekvivalentně chápat jako minimalizaci očekávané KL divergence
$$\mathbb{E}\!\left[\mathrm{KL}\!\left(q(\cdot \mid X, D_n) \,\Big\|\, \pi(\cdot \mid X, D_n)\right)\right].$$
PPD $\pi$ je tedy KL-optimálním prediktorem v $\mathcal{Q}$.

Abychom PPD aproximovali v praxi, natrénujeme model $q_\theta$, kde $\theta \in \Theta$ jsou parametry (např. váhy neuronové sítě). 
To znamená, že pro každou hodnotu $\theta$ existuje celá rodina funkcí:

$$\left\{q_{\theta} : (\mathcal{Y} \times \mathcal{X})^{n+1} \to [0,1],\; n \in \mathbb{N}\right\},$$
ale závislost na $n$ v notaci explicitně neuvádíme. 
KL-optimální parametry jsou definovány jako:
\begin{equation}
    \theta^* = \arg\max_{\theta \in \Theta} \, \mathbb{E}_{\Pi_N} \mathbb{E}_\Pi[\log q_\theta(Y \mid X, D_N)],
\label{eq:KL_optimalni_parametry}
\end{equation}
kde $\Pi_N$ je tzv. \textit{size prior}, tj. rozdělení nad velikostí trénovací množiny $N$. 
Střední hodnota přes $N$ zajišťuje, že $q_{\theta^*}$ umí napodobit celou rodinu* PPD pro různá $n$, ne pouze její $n$-tý prvek.
Pro model $q_\theta$ zpravidla neexistuje $\theta$ takové, že $q_\theta = \pi$ přesně. 
V tom případě rovnice \ref{eq:KL_optimalni_parametry} definuje \textit{KL-optimální} aproximaci $\pi$ ve třídě $\{q_\theta : \theta \in \Theta\}$, tj. nejlepší aproximaci, které je daná architektura schopna.

Dále lze ukázat následující teoretickou vlastnost.
Střední hodnota v (3) se v praxi aproximuje průměrováním přes $m$ iid datasetů, např. pomocí metody Monte Carlo.
Konkrétně generujeme datasety $(Y_j, X_j) \cup D^{(j)}$ velikosti $N_j + 1$ s $N_j \sim \Pi_N$ dle postupu \ref{generace_dat}. 
Empirická verze optimalizačního problému pak je:
\begin{equation}
    \hat{\theta} = \arg\max_{\theta \in \Theta} \sum_{j=1}^{m} \log q_\theta(Y_j \mid X_j, D^{(j)}). 
\label{optimalni_theta_empiricky}
\end{equation}

Přesnost $\hat{\theta}$ jako aproximace $\theta^*$ závisí na počtu Monte Carlo vzorků $m$.
Ale důležité je, že se dá ukázat $\hat{\theta} = \theta^* + O_p(m^{-1/2})$.
Z toho plyne, že konvergence $\hat{\theta} \to \theta^*$ je v praxi zaručena.

Nagler (2023) ukazuje, že pokud je trénovací loss NLL (negative log-likelihood, viz Fáze 8) a model má dostatečnou kapacitu, optimální řešení konverguje k Bayesovské posteriorní prediktivní distribuci. To znamená, že správně natrénovaný PFN je teoreticky ekvivalentní analytickému GP posterioru.


\subsection{Od PPD k PFN}
\label{subsec:ppd-k-pfn}
V této sekci si ukážeme podrobněji jak spolu souvisí pojmy PPD a PFN, řekneme si zajímavé vlastnosti ohledně učení PPD, která spojují tyto dva nástroje, a jako důsledek popíšeme i vlastnosti učení samotného PFN.
Uvažujme definici PPD \ref{def:PPD}.
Posterior $\Pi(\cdot \mid D_n)$ je plně určen priorem $\Pi$.
Jde o standardní Bayesovo pravidlo aplikované na nekonečnědimenzionální parametr $p$. 
Pokud jsou data $D_n$ generována z $p_0$, doufáme, že se posterior $\Pi(p \mid D_n)$ s rostoucím $n$ soustředí okolo $p_0$, a tím pádem $\pi(y \mid x, D_n) \to p_0(y \mid x)$.
Avšak v tomto okamžiku čelíme velmi netriviálnímu problému – jak zvolit správný prior, aby to konvergovalo.
Obecně existují celé články a knížky o tom, jak se vypořádat s tímto problémem, který se vyskytuje v praxi tu a tam.
Platí totiž, že je třeba, aby $\Pi$ měl dostatečně velký support a tím pokrýval dostatečně bohatou třídu funkcí.
To nám pak totiž umožní efektivní inferenci.
Např. pro případ GP s RBF kernelem bychom chtěli pokryt celý prostor všemožných tvarů té funkce, tj. aby model uměl najít správné hyperparametry v RBF kernelu.
Tím pádem, když PPD dostane sadu bodů, dokáže nám říct, ze které konkrétní konfigurace GP s RBF kernelem pochází.
Ale ani velký suport nemusí stačit.
Prior může sice $p_0$ obsahovat ve svém supportu, ale dávat příliš velkou hmotnost nepříznivým oblastem prostoru $\mathcal{P}$, čímž může konvergence posteriorů probíhat velmi pomalu nebo vůbec.
Avšak Nagler ve svém článku ukazuje, že PPD se umí učit i v situaci, kdy $p_0$ leží mimo support prioru*$\mathcal{P} = \{p : \Pi(p) > 0\}$.
To platí, za určitých podmínek.

\begin{theorem}
    Uvažujme prior $\Pi$. Nechť dále platí:
    \begin{enumerate}
        \item Existuje jedinečné $p^* \in \mathcal{P}$ takové, že
            $$p^* = \arg\min_{p \in \mathcal{P}} \, \mathrm{KL}(p \mid p_0), \qquad \mathrm{KL}(p^* \mid p_0) < \infty.$$
        \item Pro každé $\alpha \in (0, 1/2)$ existují množiny $B_1, \ldots, B_{J(\alpha)}$ takové, že
            $$\mathcal{P} \subseteq \bigcup_{j=1}^{J(\alpha)} B_j, \qquad \sup_{p, p' \in B_j} H(p, p') \leq 4(\alpha^2/2)^{1/\alpha}, \qquad \sum_{j=1}^{J(\alpha)} \Pi(B_j)^\alpha < \infty,$$
            kde $H(p, p')$ je Hellingerova vzdálenost mezi $p$ a $p'$.
    \end{enumerate}
    Pak existuje $p^* \in \mathcal{P}$ takové, že
    $$\pi(y \mid x, D_n) \xrightarrow{n \to \infty} p^*(y \mid x) \quad \text{s. j.},$$
    pro $P_0$ s.v. $(y, x)$. 
    Navíc $p^*$ je KL-optimální aproximace $p_0$ v $\mathcal{P}$:
    $$p^* = \arg\min_{p \in \mathcal{P}} \, \mathrm{KL}(p \mid p_0).$$
\label{theorem:PPDcanLearn}
\end{theorem}

První podmínka v této větě požaduje, aby v supportu prioru existovala jednoznačná KL-optimální aproximace $p^*$ pravého modelu $p_0$ a KL divergence $p^*$ vůči $p_0$ je konečná.
Druhá podmínka je technickým požadavkem na to, aby prior nepřiděloval příliš velkou hmotnost oblastem daleko od $p^*$, prior nesmí být příliš rozptýlený v nežádoucích oblastech supportu.
Sama o sobě věta tvrdí, jak se PPD učí.
Konkrétně, že se učí tím lépe, čím víc dat máme, a snaží se naučit co nejlepší aproximaci $p_0$ dosažitelnou v $\mathcal{P}$.
Pokud $p_0 \in \mathcal{P}$ (prior obsahuje pravý model): $p^* = p_0$ a PPD konverguje přímo k $p_0$.
Pokud $p_0 \notin \mathcal{P}$ (prior pravý model neobsahuje): PPD stále konverguje, ale k nejbližšímu $p^* \in \mathcal{P}$ v KL smyslu.
Logicky čím větší je $\mathcal{P}$, tím víc se přibližuje $p^*$ ke skutečnému $p_0$.
\par   

Toto tvrzení nám tedy úplně řeší problém hledání správného prioru, na který jsme upozorňovali na začátku sekce.
My totiž nepotřebujeme perfektní prior, nýbrž aby byl "dostatečně rozumný" ve smyslu předpokladů věty \ref{theorem:PPDcanLearn}, a objem dat postupně ten náš odhad prioru opraví.
Avšak PPD je ideální teoretický případ.
Pokud bychom chtěli tyto vlastnosti použít např. u PFN, je zde několik nuancí. 
PFN je totiž pouze aproximací PPD, jak ukážeme dále, trénovanou na omezených datasetech, a proto tvrzení \ref{theorem:PPDcanLearn} nelze přímo aplikovat. 
\par   

Podívejme se podrobněji na to, jak PFN aproximuje PPD.
Již jsme si definovali vztah \ref{optimalni_theta_empiricky} a teď na něj budeme nahlížet jako na optimalizační problém.
Na konečný výsledek $\hat{\theta}$ mají vliv čtyři faktory:
\begin{enumerate}
    \item datový prior $\Pi$,
    \item prior velikostí vzorků $\Pi_N$,
    \item model $q_{\theta}$,
    \item počet vzorků $m$.
\end{enumerate}

Jelikož PFN je předtrénovaný model, třídu modelů $\{q_\theta : \theta \in \Theta\}$ lze považovat za fixní vůči $m$.
Přesnost $\hat{\theta}$ jako aproximace $\theta^*$ pak plyne ze standardních výsledků empirické minimalizace vztahu: $\hat{\theta} = \theta^* + O_p(m^{-1/2})$.
Tím máme hned vyřešený jeden z faktorů.
Ostatní tři jsou mnohem zajímavější.
\par   

Pokud $\Pi$ obsahuje pouze jednoduché modely, bude i optimální PFN $q_{\theta^*}$ produkovat jednoduché funkce $(y, x)$. 
Naopak, jednoduchá architektura $\{q_\theta : \theta \in \Theta\}$ nemůže těžit ze složitého prioru $\Pi$. 
Aby PFN dobře fungoval na různorodých úlohách, musí mít dostatečnou kapacitu jak architektura $q_\theta$ jako samotná, tak i prior $\Pi$.

Prior $\Pi_N$ velikostí vzorků lze chápat jako regularizátor.
Při tréninku (odkaz) se vzorkují datasety $D^{(j)}$ s náhodnou velikostí $N_j \sim \Pi_N$. 
Definujme KL–optimální parametr pro danou velikost vzorku:
$$\theta^*_n = \arg\max_{\theta} \, \mathbb{E}_\Pi[\log q_\theta(Y \mid X, D_n)].$$
PPD $\pi(y \mid x, D_n)$, kterou se snažíme aproximovat, se mění s $n$. 
Jako funkce od $(y, x)$ roste její složitost s $n$.
Pro $n = 1$ je PPD blízká průměrnému modelu v prioru a typicky téměř konstantní.
Pro velká $n$ je stále víc a víc komplexní. 
Analogicky $\theta^*_n$ preferuje složitější modely s rostoucím $n$.
Optimalizací přes náhodné velikosti $N_j$ výsledný $\theta^*$ průměruje $\theta^*_{N_j}$. 
Distribuce $\Pi_N$ tak určuje, které velikosti datasetů jsou více signifikantní. 
To lze chápat jako regularizaci komplexity modelu.
Malé $N_j$ nutí model k jednoduchým predikcím, velké $N_j$ k složitým.
\par 

\subsection{Učení PFN In–Context}
\label{subsec:uceni-incontext}
V předchozí sekci jsme ukázali podstatnou vlastnost PPD, a to, že se se zvětšujícím se datasetem PPD přibližuje ke skutečné distribuci.
Následně jsme ukázali, že PFN je pouhou aproximací PPD, jejíž kvalita závisí na určitých faktorech, a nelze s jistotou tvrdit, že pro něj platí stejné vlastnosti jako pro PPD a nelze na něj přímo aplikovat větu \ref{theorem:PPDcanLearn}.
Ukazuje se, že PFN se pořád může chovat jako PPD a zachovat její některé klíčové vlastnosti, jako např. již zmíněné tvrzení.
Avšak v praxi to není tak jednoduché a jednoznačné.
PFN se dokáže přiblížit k $p^*(y \mid x)$, a to i při relativně malém množství dat ($n \leq 1000$).
Navíc se zachovává i to, že ani nemusíme pro tuto vlastnost najít perfektní prior.
Přesto ale prozatím nemáme žádné teoretické důvody nebo předpoklady, aby toto byla pravda.
V této sekci se pokusíme je odvodit.
\par   

V tomto okamžiku je výhodné přepnout se z bayesovského pohledu na frekventistický.
Pro libovolnou velikost datasetu $n$ při inferenci nahlížíme na předtrénovaný model $q_{\hat\theta}(y \mid x, \cdot)$ jako na nenatrénovaný frekventistický prediktor pro $p_0(y \mid x)$, což je funkce $(\mathcal{Y} \times \mathcal{X})^n \to \mathcal{P}_{\mathcal{Y}|\mathcal{X}}$, která mapuje dataset $D_n$ na podmíněnou distribuci.
Parametry $\theta$ jsou pak jen hyperparametry tohoto prediktoru, vyladěné před tréninkem. 
Prior $\Pi$ a $\Pi_N$ jsou jednoduše distribuce nad úlohami, na kterých chceme, aby prediktor fungoval dobře.
Dále rozepišme chybu predikce na bias a rozptyl:
\begin{equation}
    q_\theta(y \mid x, D_n) - p_0(y \mid x) = \underbrace{q_\theta(y \mid x, D_n) - \mathbb{E}_{D_n \sim p_0^n}[q_\theta(y \mid x, D_n)]}_{\text{rozptyl}} + \underbrace{\mathbb{E}_{D_n \sim p_0^n}[q_\theta(y \mid x, D_n)] - p_0(y \mid x)}_{\text{bias}}.
\label{Bias_Variance}
\end{equation}

Již jsme řekli, že empiricky chyba predikce klesá spolu se zvětšujícím se $n$.
Takže zkusme najít které strukturální aspekty PFN by to mohli vysvětlit.
\par 

Obecně ze struktury transformeru plyne, že jsou invariantní vůči permutacím vstupních dat (tj. symetrické), což je mimochodem jednou z výhod pro architekturu PFN dataset $D_n$ je ze své podstaty množina, ne sekvence, takže permutační invariance je přesně ta vlastnost, kterou chceme pro i.i.d. vzorky.
\begin{theorem}
    Nechť $f : (\mathcal{Y} \times \mathcal{X})^n \to \mathcal{P}_{\mathcal{Y}|\mathcal{X}}$ je libovolný prediktor. Pak existuje jeho symetrizovaná verze $\tilde{f}$ taková, že pro každé pravděpodobnostní míry $P$:
\begin{equation}
    \mathbb{E}_{D_n \sim P^n}[\tilde{f}(D_n)] = \mathbb{E}_{D_n \sim P^n}[f(D_n)], \qquad \mathrm{Var}_{D_n \sim P^n}[\tilde{f}(D_n)] \leq \mathrm{Var}_{D_n \sim P^n}[f(D_n)].
\end{equation}
\label{theorem:symetrie_transformeru}
\end{theorem}
Tato věta říká, že takové symetrické prediktory jsou optimální z hlediska MSE.
Toto je první opodstatnění faktu, proč chyba predikce pro PFN stále klesá.
\par   

Jsou tady i další strukturální vlastnosti prediktoru $q_{\theta}$, které mají vliv na kvalitu predikce.
Platí, že s větším $D_n$ klesá vliv jednotlivých vzorků na výslednou predikci. 
Formálně předpokládáme, že existují $\alpha > 0$ a $L < \infty$ takové, že pro dostatečně velká $n$ a s. v. datasety $D_n$, $D'_n$ lišící se v jediném vzorku:
\begin{equation}
    \left| q_\theta(y \mid x, D_n) - q_\theta(y \mid x, D'_n) \right| \leq L n^{-\alpha}.
\label{eq:horní_mez_prediktoru}
\end{equation}

To nám umožňuje vyslovit tvrzení, které ukazuje, proč rozptyl během učení postupně mizí.
\begin{theorem}
    Nechť platí předpoklad \ref{eq:horní_mez_prediktoru}, pak 
    \begin{equation}
        \left| q_\theta(y \mid x, D_n) - \mathbb{E}[q_\theta(y \mid x, D_n)] \right| \lesssim n^{1/2 - \alpha}
    \end{equation}
    s vysokou pravděpodobností. 
    Navíc pro $\alpha > 1/2$ platí $\lim_{n \to \infty} n^{1/2-\alpha} = 0$ a 
    \begin{equation}
        q_\theta(y \mid x, D_n) - \mathbb{E}[q_\theta(y \mid x, D_n)] \xrightarrow{n \to \infty} 0 \quad \text{almost surely.}
    \end{equation}
\end{theorem}
Jinými slovy rozptyl postupně mizí s počtem dat $n$.
Tak dostáváme, že hlavním pramenem chyby při inferenci není bias, kterého se už jen tak nezbavíme.
\par   

Ze vztahu \ref{Bias_Variance} je vidět, že bias je určen chováním posloupností $\mathbb{E}[q_\theta(y \mid x, D_n)]$.
Je logické, že bychom mohli přepodkládat (nebo očekávat), že platí i tento vztah: 
$$\mathbb{E}[q_\theta(y \mid x, D_n)] \xrightarrow{n \to \infty} \bar{q}_\theta(y \mid x),$$
kde $\bar{q}_\theta(. \mid .)$ je jakýsi zprůměrovaný prediktor, jehož konkrétní architekturu neznáme.
Ale co můžeme říct s jistotou, tak že chování biasu hodně závisí na oné architektuře, jak ukazuje Nagler ve svém článku (Nagler, 2023), může růst, klesat, anebo být celou dobu konstantní.
Co je avšak v naších sílách – vyslovit nutnou podmínku pro mizení biasu. 
Má–li prediktor $q_\theta(. \mid ., D_n)$ mizející bias na dostatečně bohaté třídě funkcí, pak nutně musí být lokální.
Lokální zde znamená, že např. k $q_\theta(y \mid x, D_n)$ by měly přispívat (asymptoticky) pouze vzorky $(Y_i, X_i) \in D_n$ s $X_i$ blízko $x$.
Tyto poznatky dává dohromady následující věta.

\begin{theorem}
    Nechť $\mathcal{P}$ je třída distribucí taková, že pro každé $p \in \mathcal{P}$:

\begin{equation}
    \mathbb{E}_{D_n \sim p^n}[q_\theta(y \mid x, D_n)] \xrightarrow{n \to \infty} p(y \mid x).
\end{equation}
    Pokud platí \ref{eq:horní_mez_prediktoru}, pak existuje posloupnost $\epsilon_n \to 0$ taková, že pro každé $\tilde{p} \in \mathcal{P}$ platí:
\begin{equation}
    \left| q_\theta(y \mid x, D_n) - q_\theta(y \mid x, \tilde{D}_n) \right| \xrightarrow{n \to \infty} 0, s.j.
\end{equation}

    kde $D_n = (Y_i, X_i)_{i=1}^n$ a $\tilde{D}_n = (Y'_i, X_i)_{i=1}^n$ s

$$Y'_i = \begin{cases} Y_i & \text{pokud } \|X_i - x\| \leq \epsilon_n, \\ \sim \tilde{p}(\cdot \mid X_i) & \text{pokud } \|X_i - x\| > \epsilon_n. \end{cases}$$
\label{theorem:lokalizace}
\end{theorem}

To znamená, že množina $\tilde{D}_n$ je poskládaná tak, že místo vzdálených od $x$ bodů obsahuje náhodný šum, ale podstatné je, že nakonec ten šum (odlehlé body) vůbec nemají vliv na výslednou predikci.
To ovšem platí pro dostatečně bohatou množinu $\mathcal{P}$.
Výsledek postrádá smysl, je-li $\mathcal{P}$ příliš malá. 
I pro bohatou $\mathcal{P}$ jde jen o nutnou, nikoliv postačující podmínku, konstantní prediktor $q_\theta = 1$ je lokální, ale jeho bias se s $n$ nemění.
Důležitý je ale důsledek pro bias-variance tradeoff: pokud bias mizí pro bohatou $\mathcal{P}$, prediktor může efektivně využít jen $\epsilon_n n = o(n)$ vzorků, takže podmínka \ref{eq:horní_mez_prediktoru} pravděpodobně neplatí pro $\alpha = 1$. 
Dále se ukáže, že lokalizace je kamenem úrazu pro transformer architekturu PFN.
\emph{Attention váhy} $a^{(h)}_j$ jsou výstupem SoftMax, tedy jsou vždy kladné pro všechny vzorky $j$, nikdy přesně nulové. 
Transformer tedy vždy váží nenulově i vzdálené vzorky. 
Přepis příznaků vzdálených bodů proto vždy predikci ovlivní, a lokalita ve smyslu věty \ref{theorem:lokalizace} není splněna. 
Transformer tak garantuje mizení rozptylu, ale nikoli mizení biasu.
Tato pozorování popíšeme podrobněji dále.

\subsection{PFN transformer}
\label{subsec:pfn-transformer}
Jelikož tato práce je věnovaná pohledu na PFN skrz architekturu transformer (Vaswani et al., 2017), dává smysl podívat se na statistické vlastnosti právě této architektury.
Uvažujme tedy opět  pro jednoduchost transformer s jednou vrstvou. 
Dataset $D_n = \{V_i\}_{i=1}^n$, kde $V_i = (Y_i, X_i) \in \{0,1\} \times \mathbb{R}^d$, a testovací vektor $v = (0, x)$. 
Síť je definována následujícími operacemi:

\begin{gather*}
a_j^{(h)} = \mathrm{SoftMax}\left(v^\top W_q^{(h)}V_1, \dots, v^\top W_q^{(h)}V_n\right)_j, \\[2ex]
u' = \sum_{h=1}^H \sum_{j=1}^n a_j^{(h)} W_v^{(h)} V_j, \\[2ex]
u = \mathrm{LayerNorm}(v + u'; \gamma), \\[2ex]
z' = W_{r,2}\mathrm{ReLU}(W_{r,1}u; \gamma), \\[2ex]
z = \mathrm{LayerNorm}(u + z'; \gamma), \\[2ex]
q_\theta(\cdot \mid x, D_n) = \mathrm{SoftMax}(W_o z),
\end{gather*}

kde matice vah jsou $W^{(h)}_q, W^{(h)}_v \in \mathbb{R}^{(d+1)\times(d+1)}$, $W_{r,1}, W_{r,2}^\top \in \mathbb{R}^{m\times(d+1)}$, $W_o \in \mathbb{R}^{|\mathcal{Y}|\times(d+1)}$. 
V tomto případě parametr $\theta$, který jsme všude uváděli v kontextu PFN v předchozích sekcích, sbírá všechny tyto matice. 
To znamená, že $\theta = \{ W_q^{(1)}, \dots, W_q^{(H)}, W_v^{(1)}, \dots, W_v^{(H)}, W_{r,1}, W_{r,2}, W_o, \gamma \}$
Operace SoftMax, LayerNorm a ReLU jsou definovány jako:

\begin{gather*}
\mathrm{SoftMax}(v) = \frac{\exp(v)}{\sum_j \exp(v_j)}, \qquad \mathrm{LayerNorm}(v;\,\gamma) = \gamma_1 \frac{v - \mathrm{avg}(v)}{\|v - \mathrm{avg}(v)\|_2 + |\gamma_2|} + \gamma_3,
\mathrm{ReLU}(v) = \max(0, v),
\end{gather*}

kde $\exp$ a $\max$ působí po složkách.
První dvě rovnice definují \emph{attention mechanismus} s $H$ hlavami (Vaswani et al., 2017). 
Ze součtu $a^{(h)}_1 + \cdots + a^{(h)}_n = 1$ plyne, že \emph{attention váhy} tvoří pravděpodobnostní distribuci přes trénovací vzorky.
Myšlenka je, že v každé hlavě \emph{attention váhy} $a^{(h)}_j$ zdůrazňují konkrétní vzorky $V_j \in D_n$, konkrétně ty, které jsou "podobné" testovacímu vektoru $v$ ve smyslu měřeném skalárním součinem $v^\top W^{(h)}_q V_j$. 
Každá \emph{attention hlava} dovoluje jinou definici podobnosti prostřednictvím matice $W^{(h)}_q$.
Intuitivně to funguje takto: transformer si při predikci pro $x$ vyhledá v trénovacích datech vzorky, které považuje za relevantní, a jejich labely agreguje. Různé hlavy mohou sledovat různé aspekty podobnosti, jedna hlava třeba srovnává hodnoty $X_i$ s $x$, jiná sleduje jiný rys vstupních vektorů.
\par   

Nás by také zajímalo, jaké vlastnosti má rozptyl a bias pro tuto architekturu, protože v předchozí kapitole jsme poznamenávali, že konkretní vlastnosti těchto charakteristik se mění v závislosti na volbě konkrétní architektury.
\begin{theorem}
    Pro $\mathcal{X} = \{x : \|x\|_2 \leq K\}$ a omezené matice vah ($\|W^{(h)}_q\|_2, \|W^{(h)}_v\|_2, \|W_{r,1}\|_2, \|W_{r,2}\|_2, \|W_o\|_2 < \infty$) platí
    \begin{equation}
            \left|q_\theta(y \mid x, D_n) - \mathbb{E}[q_\theta(y \mid x, D_n)]\right| \lesssim n^{-1/2}
    \end{equation}
s vysokou pravděpodobností.
\label{theorem:Rozptyl_transformeru}
\end{theorem}

Velmi podobnou větu jsme již viděli v předchozí sekci.
Dokonce tvrdí skoro stejné: rozptyl mizí nezávisle na parametru $\theta$.
Jde o strukturální vlastnost architektury, nikoliv výsledek tréninku. 
Důvodem je, že \emph{attention mechanismus} nutně přiděluje každému vzorku váhu řádu $1/n$, takže vliv libovolného jednotlivého vzorku klesá s $n$.
V případě biasu jsme již říkali, že ta situace je horší. 
Bias v případě transformeru totiž silně závisí na volbě $\theta$.
Označme $\bar{q}_\theta(y \mid x)$ limitní hodnotu predikce pro $n \to \infty$.
\begin{theorem}
    Za předpokladů platnosti věty \ref{theorem:Rozptyl_transformeru} plyne
    \begin{equation}
            \mathbb{E}[q_\theta(\cdot \mid x, D_n)] \xrightarrow{n \to \infty} \bar{q}_\theta(\cdot \mid x),
    \end{equation}
    kde $\bar{q}_\theta(\cdot \mid x)$ je definováno stejně jako $q_\theta(\cdot \mid x, D_n)$, ale s $u'$ nahrazeným výrazem $\bar{u}' = \sum_{h=1}^{H} W^{(h)}_v \mathbb{E}_{V \sim g_h}[V],$
    kde $g_h(s) = \frac{\exp(v^\top W^{(h)}_q s)}{\mathbb{E}_{V \sim p_0}[\exp(v^\top W^{(h)}_q V)]} \cdot p_0(s).$
\end{theorem}

To znamená, že v definici transformeru výše jsme nahradili výraz $u'$ tím novým $\hat{u'}$.
Navíc v definici $\hat{u'}$ se objevuje tzv. exponentially tilted function (česky exponenciálně vychýlená funkce).
V tomto kontextu se zavádí jako asymptotická variace \emph{attention mechanismu}.
Tato funkce $g_h(s)$ zvyšuje pravděpodobnost těm hodnotám $s$, které se nejvíce podobají vektoru $v$ ve smyslu $v^\top W^{(h)}_q s$, a snižuje pro zbylé body.
Jinými slovy ten složitý vzoreček na obrázku ukazuje teoretickou situaci, co by se stalo, kdyby měl transformer k dispozici nekonečně mnoho vzorků. 
Už nerozděluje váhy mezi jednotlivými datovými body, ale vezme celou spojitou distribuci dat a "vychýlí" ji celou směrem k té informaci, na kterou zrovna potřebuje upřít pozornost.
Více formálně to znamená, že každá \emph{attention hlava} $h$ tak hodnotí určitý aspekt neznámé distribuce $p_0$ (charakterizovaný maticí $W^{(h)}_q$). 
Pokud jsou matice $(W^{(h)}_q)_{h=1}^H$ dobře nastaveny, tyto jednotlivé pohledy rozlišují různé hodnoty příznaků.
\par   

Avšak $g_h$ lokalizuje prediktor pouze do jisté míry.
Ačkoli zvyšuje váhu vzorků podobných $v$, ale ne ve smyslu věty \ref{theorem:lokalizace}.
Vzorky vzdálené od $x$ stále nenulově přispívají, protože SoftMax nikdy nevytváří přesně nulové váhy. 
Přepsání příznaků vzdálených vzorků proto vždy predikci změní, a tak lokalita podle \ref{theorem:lokalizace} není splněna. 
Proto plyne, že bias transformeru obecně nemůže úplně zmizet.
Limitní bias $\bar{q}_\theta(y \mid x) - p_0(y \mid x)$ přesto může být malý, pokud síť za \emph{attention vrstvou} dokáže ze součtu aspektových souhrnů $W^{(h)}_v \mathbb{E}_{V \sim g_h}[V]$ zkonstruovat dobrou aproximaci $p_0$. 
Relevance jednotlivých aspektů závisí na pravé distribuci $p_0$.
Méně relevantní aspekty mohou přispívat méně. 
Na malých vzorcích ($n = 1$) jsou všechny \emph{attention váhy} $a^{(h)}_j = 1$, takže všechny aspekty přispívají stejnou měrou. 
To naznačuje, že bias může klesat v rozsahu velikostí, na které byl $\theta$ natrénován, nikoliv však za jejich hranicí, kde není důvod očekávat pokles biasu.
Klíčovou roli hraje přítomnost více \emph{attention hlav} ($H > 1$), spolupráce více hlav může efektivně napodobit mechanismus průměrování modelů.
Dalším klíčovým nástrojem je větší počet \emph{attention vrstev}.
Navíc empirické testy ukazují, že na výslednou kvalitu predikce (a zmenšování biasu) hraje neposlední roli samotný trénink a jeho nastavení.
Všechny tyto manipulační součástky tak dávají velký prostor pro manévr a různé jejich kombinace či nastavení mohou vést k velmi signifikantnímu poklesu biasu.
\par

Dále lze různě upravovat strukturu PFN během inference či tréninku pro zlepšení lokality.
Naším účelem je tedy nějakým způsobem donutit transformer se dívat na jednotlivé body co nejvíc lokálně a neuvažovat pro bod $x$ příliš vzdálené hodnoty.
Jedním z přímočarých způsobů je \textit{post–hoc lokalizace}, kterou navrhuje Nagler.
Pro predikci labelu v testovacím bodě $x$:
\begin{enumerate}
    \item Sestrojí se redukovaná trénovací množina $D_n(x)$ odstraněním všech vzorků kromě $k_n$ nejbližších sousedů bodu $x$ z $D_n$.
    \item Predikuje se label maximalizující $q_\theta(\cdot \mid x, D_n(x))$.
\end{enumerate}

Omezení na okolí $x$ je ekvivalentní "roztažení" cílové funkce $p(y \mid \cdot)$ v okolí $x$, lokálně hladká funkce se chová přibližně jako konstantní. 
Konstantní funkce jsou pro $q_\theta$ snazší k aproximaci. 
Lokalizace tedy zlepšuje bias za cenu mírného zvýšení rozptylu, prediktor pracuje s menší efektivní velikostí datasetu $k_n \ll n$.
